% Draft paper based on the KAIA LaTeX template.

\documentclass[nouppercase]{kaia}
\usepackage{amsmath}
\usepackage{amssymb}
\usepackage{graphicx}

\title{Selective Temporal Mask Expansion for Reducing Residue and Over-Removal in Video Object Removal}

\author{First Author}{a}
\author{Second Author}{a}
\affiliation{Department of Artificial Intelligence, Institution Name, City, Republic of Korea, \{first.author, second.author\}@example.com}{a}

\begin{document}

\maketitle
\begin{abstract}
Video object removal removes a target object from every frame and fills the missing region with temporally coherent background content. Existing video inpainting models can produce plausible per-frame results, but object residue and boundary flicker often remain when the input mask is temporally inconsistent or too conservative. Boundary-only mask expansion can reduce local boundary artifacts, but it often fails to remove residual object pixels. In contrast, using a full temporal union of masks removes more residue but frequently changes background regions that should remain untouched. This paper proposes selective temporal mask expansion, a training-free and model-agnostic preprocessing strategy that selects reliable extra regions from a temporal mask union. The method starts from a boundary-expanded mask and adds extra temporal-union pixels only when they lie near the object boundary and have sufficient temporal occupancy. This produces a controllable trade-off between residue removal and over-removal suppression. Experiments on 100 AI-Hub video clips and 30 DAVIS 2017 validation clips show that the proposed method reduces residue relative to boundary-only expansion while consistently reducing outside-region changes relative to full temporal union. On the AI-Hub 100-clip set, the main variant improves residue over boundary-only in 98/100 clips and reduces outside changes relative to temporal union in 97/100 clips. Parameter ablations over radius and occupancy threshold show that the main setting is a balanced operating point, and threshold sensitivity tests at RGB difference thresholds of 5, 10, and 20 show the same trade-off pattern. These results support selective temporal expansion as a practical mask refinement strategy for existing video inpainting pipelines.
\end{abstract}

\begin{keywords}
Video Inpainting, Object Removal, Mask Refinement, Temporal Consistency, Video Segmentation
\end{keywords}

\section{Introduction}
Video object removal is a practical video inpainting task in which an unwanted object is removed from a sequence and the missing region is reconstructed from spatial and temporal context. Recent methods such as ProPainter~\cite{zhou2023propainter} have improved propagation and long-range video reasoning, but the final quality still depends strongly on the input masks. If the masks are too small or temporally unstable, object residue can remain near the boundary. If the masks are expanded too aggressively, the inpainting model may unnecessarily alter background content.

This trade-off is especially visible around object boundaries. Boundary-only smoothing or expansion is a conservative choice: it modifies a narrow band around the object and can reduce small boundary artifacts. However, it does not address residue that appears away from the immediate boundary or appears intermittently across frames. A full temporal union of the masks is the opposite extreme. It covers pixels that the object occupies at any time in the sequence and therefore removes more residue, but it can create a large mask that includes valid background regions in many frames. This often causes over-removal artifacts and unnecessary background changes.

This paper studies the following question: can we obtain much of the residue-removal benefit of temporal union without accepting its full over-removal cost? We propose selective temporal mask expansion. Instead of using the full temporal union, the method uses the temporal union as a pool of candidate extra pixels and selects only pixels that are close to the object boundary and sufficiently supported over time. The selected extra region is then added to a boundary-expanded base mask and passed to an existing video inpainting model. The method is training-free, model-agnostic, and can be applied as a plug-in mask refinement stage before running ProPainter or similar inpainting models.

Our framing differs from depth- or flow-only temporal refinement. We do not claim that a single cue explains all flicker or residue. Instead, we treat mask refinement as a controlled trade-off between two failure modes: residual object pixels and over-removal of background. The contribution is an occupancy-guided temporal selection rule that can be evaluated independently of network training or architecture changes.

The main contributions are as follows:
\begin{itemize}
    \item We analyze the trade-off between boundary-only expansion and full temporal union for video object removal masks.
    \item We propose a selective temporal mask expansion strategy that adds temporally supported extra mask regions while suppressing unnecessary outside changes.
    \item We evaluate the method on 100 AI-Hub video clips and 30 DAVIS 2017 validation clips using residue, outside-change, extra-mask, and boundary temporal-error metrics.
    \item We provide clip-level win counts and qualitative success/failure cases to show both the strengths and limitations of the method.
\end{itemize}

\section{Related Work}
\subsection{Video Inpainting and Object Removal}
Video inpainting reconstructs missing regions across a sequence while maintaining spatial plausibility and temporal coherence. Earlier deep methods jointly modeled temporal structure and spatial details~\cite{wang2018video}, while flow-guided methods propagated information using completed optical flow~\cite{xu2019deep}. ProPainter~\cite{zhou2023propainter} improves video inpainting using dual-domain propagation and a mask-guided sparse video transformer. Although these models are strong, their output quality can still be limited by imperfect masks.

\subsection{Mask Quality in Video Object Removal}
Object removal is often treated as an inpainting problem after a mask has already been produced. In practice, however, the mask is a key part of the system. Conservative masks may preserve background but leave object residue. Aggressive masks may remove residue but change regions that do not need inpainting. This paper focuses on this mask-level trade-off rather than modifying the inpainting network.

\subsection{Temporal Propagation}
Optical flow models such as RAFT~\cite{teed2020raft} are widely used for video alignment and temporal propagation. Flow-based refinement can reduce flicker, but it may be unreliable near occlusion boundaries, fast motion, and thin structures. Our method does not require training a flow-based refinement module. Instead, it uses temporal mask occupancy as a simple signal for selecting candidate regions from a temporal union.

\section{Method}
\subsection{Problem Setup}
Let $I_t$ be the input frame at time $t$ and $M_t$ be the original binary object mask. A video inpainting model receives the sequence of frames and masks and produces an output $O_t$. The goal is to improve the input masks so that the final output reduces object residue while avoiding unnecessary changes outside the target object region.

We compare five mask strategies:
\begin{itemize}
    \item \textbf{Original}: use the raw mask $M_t$.
    \item \textbf{Boundary-only}: expand only a narrow object boundary band.
    \item \textbf{Temporal union}: use the full temporal union of masks.
    \item \textbf{Ours r10 t0.10}: selective expansion with radius $r=10$ and occupancy threshold $\tau=0.10$.
    \item \textbf{Ours r10 t0.15}: a conservative selective expansion with $r=10$ and $\tau=0.15$.
\end{itemize}

\subsection{Temporal Union}
The temporal union mask is defined as
\begin{equation}
    U = \bigcup_{t=1}^{T} M_t.
\end{equation}
Using $U$ at every frame can cover object locations across the entire sequence. This is effective for removing residue caused by mask jitter or motion, but it can also be too aggressive. If the object moves across a large area, $U$ may include background regions in frames where the object is absent.

\subsection{Boundary-Expanded Base Mask}
We first build a boundary-expanded base mask. Let $\mathcal{D}_r(\cdot)$ and $\mathcal{E}_r(\cdot)$ denote dilation and erosion with radius $r$. The boundary band is
\begin{equation}
    B_t^r = \mathcal{D}_r(M_t) - \mathcal{E}_r(M_t).
\end{equation}
The boundary-only baseline adds temporal-union candidates only in a local boundary band:
\begin{equation}
    M_t^{base} = M_t \cup \left((U - M_t) \cap B_t^r \right).
\end{equation}
This mask is conservative and mainly targets boundary flicker.

\subsection{Selective Temporal Expansion}
The proposed method uses the temporal union as a candidate pool but selects only reliable extra pixels. We compute a temporal occupancy map from the mask sequence:
\begin{equation}
    A(x) = \frac{1}{T}\sum_{t=1}^{T} M_t(x).
\end{equation}
Pixels with larger $A(x)$ are occupied by the object in more frames and are more likely to represent consistent object support rather than a one-frame mask error.

For each frame, the extra candidate region is
\begin{equation}
    E_t = U - M_t^{base}.
\end{equation}
The selected extra region is
\begin{equation}
    S_t = E_t \cap \mathcal{D}_r(M_t) \cap \{x \mid A(x) > \tau\}.
\end{equation}
The final mask is
\begin{equation}
    M_t^{ours} = M_t^{base} \cup S_t.
\end{equation}
The radius $r$ controls how far the method can expand from the current object boundary, and the threshold $\tau$ controls how conservative the temporal support must be. We report $r=10,\tau=0.10$ as the main setting and $r=10,\tau=0.15$ as a conservative variant.

\section{Evaluation Metrics}
We evaluate mask strategies through the final ProPainter output rather than mask geometry alone. The following metrics are computed per frame and averaged.

\subsection{Boundary Temporal Error}
Boundary temporal error measures flicker near the object boundary. We estimate optical flow between adjacent original frames, warp the previous output to the current frame, and compute the temporal difference inside the boundary band:
\begin{equation}
    \text{BTE} =
    \frac{1}{T-1}\sum_{t=2}^{T}\frac{1}{|B_t|}
    \sum_{x \in B_t} \left|O_t(x)-\mathcal{W}(O_{t-1})(x)\right|.
\end{equation}
Lower values indicate better temporal stability near the boundary.

\subsection{Outside Changed}
Outside changed measures the fraction of non-object pixels whose output differs from the original input by more than a threshold:
\begin{equation}
    \text{Outside} =
    \frac{|\{x \notin M_t : |O_t(x)-I_t(x)| > \delta\}|}{|\{x : x \notin M_t\}|}.
\end{equation}
This is a proxy for over-removal. We use $\delta=10$ in RGB intensity units.

\subsection{Extra Mask Ratio and Residue Proxy}
The extra mask ratio measures how much additional mask area is introduced relative to the original mask:
\begin{equation}
    \text{Extra} = \frac{|M_t^{variant} - M_t|}{|M_t|}.
\end{equation}
For residue, we use a simple output-input similarity proxy inside the original mask:
\begin{equation}
    \text{Residue} =
    \frac{|\{x \in M_t : |O_t(x)-I_t(x)| \leq \delta\}|}{|M_t|}.
\end{equation}
Lower values suggest that fewer original object pixels remain visually similar to the input.

\section{Experiments}
\subsection{Datasets}
The main experiment uses 100 clips from an AI-Hub video inpainting dataset. These clips contain diverse object categories, indoor/outdoor scenes, and temporally varying masks. We also evaluate on the DAVIS 2017 validation set~\cite{pont2017davis}, using 30 clips as an external validation set. DAVIS is not originally designed for object removal, but its high-quality segmentation masks make it useful for testing whether the same mask-expansion trade-off appears in a different dataset.

\subsection{Implementation Details}
All methods use ProPainter as the inpainting backend. We generate a separate inpainting result for each mask variant. For AI-Hub, the 100-clip set contains 3973 evaluated frames. For DAVIS, all 30 validation clips are processed, for a total of 1999 frames. Boundary temporal error is computed using dense optical flow estimated from adjacent original frames. Unless otherwise stated, outside-change and residue metrics use an RGB difference threshold of $\delta=10$. The proposed main setting is $r=10,\tau=0.10$, and the conservative variant is $r=10,\tau=0.15$.

\section{Results}
\subsection{Main Quantitative Results}
Table~\ref{tab:driving} shows results on the AI-Hub 100-clip set. Boundary-only slightly improves boundary temporal error but barely reduces residue. Temporal union strongly reduces residue but greatly increases outside changed and extra mask area. The proposed method lies between these extremes: it reduces residue substantially compared with boundary-only while keeping outside changes much lower than temporal union.

\begin{table}[h]
\begin{center}
\resizebox{\linewidth}{!}{%
\begin{tabular}{|l|c|c|c|c|}
\hline
Method & BTE $\downarrow$ & Outside $\downarrow$ & Extra $\downarrow$ & Residue $\downarrow$ \\
\hline\hline
Original & 0.017954 & 0.060028 & 0.0000 & 0.4231 \\
Boundary-only & 0.017147 & 0.061762 & 0.1041 & 0.4154 \\
Temporal union & 0.015304 & 0.276760 & 3.6329 & 0.1465 \\
Ours r10 t0.10 & 0.014490 & 0.177325 & 1.7081 & 0.2077 \\
Ours r10 t0.15 & 0.014493 & 0.148977 & 1.2078 & 0.2332 \\
\hline
\end{tabular}
}
\end{center}
\caption{Results on the AI-Hub 100-clip set. The proposed variants provide a trade-off between residue removal and over-removal suppression.}
\label{tab:driving}
\end{table}

Table~\ref{tab:davis} shows the same trend on DAVIS 2017 validation clips. Temporal union again produces the lowest residue and boundary temporal error, but at the cost of very large outside changes. The proposed variants reduce outside changes substantially relative to temporal union while retaining part of the residue-removal benefit.

\begin{table}[h]
\begin{center}
\resizebox{\linewidth}{!}{%
\begin{tabular}{|l|c|c|c|c|}
\hline
Method & BTE $\downarrow$ & Outside $\downarrow$ & Extra $\downarrow$ & Residue $\downarrow$ \\
\hline\hline
Original & 0.042671 & 0.008632 & 0.0000 & 0.0655 \\
Boundary-only & 0.041052 & 0.014502 & 0.3402 & 0.0622 \\
Temporal union & 0.034287 & 0.162685 & 4.6084 & 0.0575 \\
Ours r10 t0.10 & 0.035826 & 0.092581 & 2.7407 & 0.0609 \\
Ours r10 t0.15 & 0.036236 & 0.076194 & 2.2411 & 0.0606 \\
\hline
\end{tabular}
}
\end{center}
\caption{External validation on DAVIS 2017 val clips. DAVIS follows the same trade-off pattern as the AI-Hub clips.}
\label{tab:davis}
\end{table}

\subsection{Parameter Ablation}
To verify that the reported setting is not a cherry-picked point, Table~\ref{tab:param_ablation} evaluates radius and occupancy-threshold sensitivity on the AI-Hub 100-clip set. With $\tau=0.10$ fixed, increasing $r$ from 5 to 15 changes the trade-off only mildly: outside changes remain near 0.177, while residue decreases from 0.2093 to 0.2061 and BTE decreases from 0.014526 to 0.014421. This indicates that the method is not overly sensitive to a single radius choice around $r=10$. With $r=10$ fixed, $\tau$ controls the expected residue/outside trade-off. A permissive threshold ($\tau=0.05$) removes more residue but increases outside changes and mask area, while conservative thresholds ($\tau=0.15,0.20$) reduce outside changes but leave more residue. The main setting $r=10,\tau=0.10$ is therefore used as a balanced point, and $r=10,\tau=0.15$ is reported as a conservative operating point.

\begin{table}[h]
\begin{center}
\begin{tabular}{|c|c|c|c|c|}
\hline
$r$ & $\tau$ & Outside $\downarrow$ & Residue $\downarrow$ & BTE $\downarrow$ \\
\hline\hline
5 & 0.10 & 0.176974 & 0.2093 & 0.014526 \\
10 & 0.10 & 0.177325 & 0.2077 & 0.014490 \\
15 & 0.10 & 0.177788 & 0.2061 & 0.014421 \\
\hline
10 & 0.05 & 0.220388 & 0.1778 & 0.014608 \\
10 & 0.10 & 0.177325 & 0.2077 & 0.014490 \\
10 & 0.15 & 0.148977 & 0.2332 & 0.014493 \\
10 & 0.20 & 0.131361 & 0.2520 & 0.014811 \\
\hline
\end{tabular}
\end{center}
\caption{Parameter ablation on the AI-Hub 100-clip set. The first block varies radius with $\tau=0.10$ fixed, and the second block varies occupancy threshold with $r=10$ fixed.}
\label{tab:param_ablation}
\end{table}

Figure~\ref{fig:pareto} visualizes the same result as an outside-residue trade-off. Boundary-only stays close to the original outside-change level but leaves high residue, whereas temporal union reaches low residue by moving far to the high-outside-change region. The proposed variants occupy the middle region of the plot. The radius sweep forms a tight cluster, which supports the interpretation that radius has a limited effect around the chosen range. In contrast, the occupancy-threshold sweep moves along a clear trade-off curve: lowering $\tau$ prioritizes residue removal, while raising $\tau$ prioritizes background preservation. This behavior motivates $r=10,\tau=0.10$ as the residue-oriented main setting and $r=10,\tau=0.15$ as the conservative setting.

\begin{figure}[h]
\begin{center}
\includegraphics[width=\linewidth]{../../experiments/paper_visualizations/pareto/aihub100_outside_residue_pareto.png}
\end{center}
\caption{Pareto-style outside-residue trade-off on the AI-Hub 100-clip set. Lower-left is better. The occupancy threshold controls the main trade-off direction, while the radius sweep remains comparatively stable.}
\label{fig:pareto}
\end{figure}

\subsection{Clip-Level Win Counts}
Mean scores can hide whether the behavior is consistent across clips. Table~\ref{tab:wincount} reports clip-level win counts. On the AI-Hub 100-clip set, the main method improves residue over boundary-only in 98/100 clips and reduces outside changes relative to temporal union in 97/100 clips. The conservative variant improves residue in 97/100 clips and reduces outside changes in 98/100 clips. On DAVIS, the method also reduces outside changes in all 30 clips. These counts show that the method is not driven by only a few favorable sequences.

\begin{table}[h]
\begin{center}
\resizebox{\linewidth}{!}{%
\begin{tabular}{|l|l|c|c|c|c|c|}
\hline
Dataset & Method & Residue & Outside & Both & BTE & Gap \\
 & & $<$ B-only & $<$ Union &  & $<$ Orig. & $>$0.10 \\
\hline\hline
AI-Hub 100 & Ours r10 t0.10 & 98/100 & 97/100 & 96/100 & 75/100 & 20/100 \\
AI-Hub 100 & Ours r10 t0.15 & 97/100 & 98/100 & 96/100 & 83/100 & 36/100 \\
DAVIS & Ours r10 t0.10 & 21/30 & 30/30 & 21/30 & 26/30 & 0/30 \\
DAVIS & Ours r10 t0.15 & 20/30 & 30/30 & 20/30 & 27/30 & 0/30 \\
\hline
\end{tabular}
}
\end{center}
\caption{Clip-level win counts. ``Gap'' counts clips where the residue gap against temporal union is larger than 0.10, which indicates a limitation of selective expansion.}
\label{tab:wincount}
\end{table}

\subsection{Threshold Sensitivity and Clip-Level Confidence}
Table~\ref{tab:threshold} evaluates the AI-Hub 100-clip set with RGB difference thresholds $\delta=5,10,20$. Across all thresholds, the proposed variants keep residue much lower than Boundary-only while keeping outside changes lower than Temporal union. This indicates that the observed trade-off is not an artifact of a single threshold choice.

\begin{table}[h]
\begin{center}
\begin{tabular}{|c|l|c|c|}
\hline
$\delta$ & Method & Outside $\downarrow$ & Residue $\downarrow$ \\
\hline\hline
5 & Boundary-only & 0.157211 & 0.2912 \\
5 & Temporal union & 0.367517 & 0.0757 \\
5 & Ours r10 t0.10 & 0.273788 & 0.1211 \\
5 & Ours r10 t0.15 & 0.245198 & 0.1409 \\
\hline
10 & Boundary-only & 0.061762 & 0.4154 \\
10 & Temporal union & 0.276760 & 0.1465 \\
10 & Ours r10 t0.10 & 0.177325 & 0.2077 \\
10 & Ours r10 t0.15 & 0.148977 & 0.2332 \\
\hline
20 & Boundary-only & 0.017889 & 0.5447 \\
20 & Temporal union & 0.204636 & 0.2647 \\
20 & Ours r10 t0.10 & 0.113262 & 0.3352 \\
20 & Ours r10 t0.15 & 0.088747 & 0.3639 \\
\hline
\end{tabular}
\end{center}
\caption{Threshold sensitivity on the AI-Hub 100-clip set. The same residue/outside-change trade-off holds for all tested thresholds.}
\label{tab:threshold}
\end{table}

We also compute clip-level bootstrap confidence intervals on the AI-Hub 100-clip set. For Ours r10 t0.10, the mean residue gain over Boundary-only is 0.2066 with a 95\% confidence interval of [0.1829, 0.2307], and the mean outside-change reduction relative to Temporal union is 0.0989 with a 95\% confidence interval of [0.0767, 0.1228]. For Ours r10 t0.15, the corresponding values are 0.1812 [0.1593, 0.2035] and 0.1270 [0.1013, 0.1550].

\subsection{Computational Cost and Practicality}
Table~\ref{tab:cost} summarizes the practical cost of each mask strategy. All methods use the same ProPainter backend and require no additional training. The proposed method uses only the binary mask sequence: temporal union, temporal occupancy, and boundary support are computed with lightweight CPU image operations. Optical flow is used only for the boundary temporal error metric, not for mask generation. In the experiments, each mask variant is inpainted separately so that the trade-off can be measured, but deployment requires only one ProPainter run after choosing the desired operating point.

\begin{table}[h]
\begin{center}
\resizebox{\linewidth}{!}{%
\begin{tabular}{|l|c|c|c|l|}
\hline
Method & Training & Model change & Deploy runs & Extra preprocessing \\
\hline\hline
Original mask & No & No & 1 & None \\
Boundary-only & No & No & 1 & Boundary dilation from binary masks \\
Temporal union & No & No & 1 & Temporal max over binary masks \\
Ours & No & No & 1 & Union, occupancy, and boundary filtering \\
\hline
\end{tabular}
}
\end{center}
\caption{Computational cost and practicality. The proposed method is a training-free, model-agnostic preprocessing step and does not require depth, optical flow, or model modification at inference time.}
\label{tab:cost}
\end{table}

\subsection{Qualitative Results}
We select representative qualitative cases using metric-based clip rankings computed from the AI-Hub 100-clip frame metrics. The ranking includes balanced success, temporal-union over-removal, residue limitation relative to temporal union, and over-removal limitation relative to boundary-only. The temporal-union failure case in Figure~\ref{fig:cases} is ranked third by the temporal-union over-removal score, and the limitation case is ranked within the top 15 by the proposed method's over-removal score. Figure~\ref{fig:cases} summarizes representative qualitative cases. The driving success case shows that boundary-only can leave residue, while temporal union changes more background content. The DAVIS success case confirms that the same trend appears on a different dataset. The temporal-union failure case shows a large outside-region change caused by an overly broad union mask. The final case shows a limitation of the proposed method: when the selected extra region is still too broad, it can also increase outside changes.

\begin{figure*}[t]
\begin{center}
\includegraphics[width=0.95\linewidth]{../../experiments/paper_visualizations/case_panels/01_success_driving_driving_I-210715_I03012_W06_frame_0016_crop.png}
\includegraphics[width=0.95\linewidth]{../../experiments/paper_visualizations/case_panels/02_success_davis_davis_motocross-jump_00028_crop.png}
\includegraphics[width=0.95\linewidth]{../../experiments/paper_visualizations/case_panels/03_temporal_union_failure_driving_I-210618_I01001_W01_frame_0011_crop.png}
\includegraphics[width=0.95\linewidth]{../../experiments/paper_visualizations/case_panels/04_ours_failure_driving_I-210720_O12052_T04_frame_0001_crop.png}
\end{center}
\caption{Representative cases. From top to bottom: success on AI-Hub clips, success on DAVIS, temporal-union over-removal, and a limitation case of the proposed method. The top row of each panel compares outputs and the bottom row visualizes change heatmaps against the input frame.}
\label{fig:cases}
\end{figure*}

\section{Ablation and Discussion}
The ablation compares the main mask families and also sweeps the two parameters of selective temporal expansion. Original masks represent the raw input, boundary-only expansion tests conservative boundary smoothing, and temporal union tests full temporal expansion. The parameter sweep shows that $r=10,\tau=0.10$ is a balanced operating point rather than an isolated favorable setting. The conservative setting, $r=10,\tau=0.15$, reduces outside changes further at the cost of weaker residue removal. The threshold-sensitivity experiment further shows that this ordering is stable under different output-input difference thresholds.

The results suggest that selective expansion should be understood as a trade-off method rather than a universal replacement for temporal union. Temporal union can remove residue more aggressively, and in some clips the proposed method has a clear residue gap relative to union. However, temporal union also changes background regions much more often. The proposed method is useful when over-removal is a major concern and when a moderate residue trade-off is acceptable.

\subsection{Preliminary Soft-Confidence Extension}
As a preliminary extension, we also tested whether temporal occupancy can be used as a soft confidence signal rather than only as a hard mask-selection criterion. We generated confidence maps from temporal occupancy, boundary distance, and their product, and used them as post-hoc output blending gates between the ProPainter output and the original frame. This strongly reduced outside-region changes: for example, temporal-union output with occupancy-distance blending reduced Outside from 0.276760 to 0.008502 while keeping the residue proxy at 0.1465. However, the same variant increased Boundary TE from 0.015304 to 0.043342, indicating temporal seams caused by mixing original and inpainted pixels after inference. Therefore, we do not treat output blending as the main method in this paper. Instead, soft confidence guidance is left as future work, where temporally smoothed gates or model-level attention bias may use the confidence map without introducing post-hoc blending artifacts.

\section{Limitations}
The method has several limitations. First, temporal occupancy is a mask-level cue and does not understand scene geometry. If the object covers a broad area over time, the selected extra region can still become too large. Second, the residue metric is a proxy based on output-input similarity inside the original mask; it may not perfectly match human perception. Third, DAVIS is used as an external validation dataset rather than a native object-removal benchmark. Finally, the experimental comparison requires rerunning the inpainting model for each mask variant, although deployment uses only the selected mask variant.

\section{Conclusion}
This paper proposed selective temporal mask expansion for video object removal. The method selects temporally supported extra mask regions from a full temporal union and adds them to a boundary-expanded base mask. Experiments on the AI-Hub 100-clip set and DAVIS 2017 validation clips show that the method consistently reduces outside-region changes relative to temporal union while improving residue over boundary-only expansion in most clips. The method is training-free, model-agnostic, and compatible with existing video inpainting models, making it a practical mask refinement strategy for reducing the trade-off between object residue and over-removal.

\bibliography{depth_aware_temporal_refinement_refs}

\newpage
\newpage

\onecolumn
\begin{center}
    \section*{Summary of this paper}
    \setcounter{subsection}{0}
\end{center}
\noindent\fbox{
    \parbox{\textwidth}{
        \begin{quote}
        \subsection{Problem Setup}
        Video object removal quality depends strongly on mask quality. Conservative masks preserve background but can leave object residue, while full temporal union masks remove more residue but can over-edit background regions.
        \vspace{0.5\baselineskip}

        \subsection{Novelty}
        The paper proposes selective temporal mask expansion, a training-free mask refinement method that uses temporal union as a candidate pool but selects only boundary-near and temporally supported extra regions.
        \vspace{0.5\baselineskip}

        \subsection{Algorithms}
        The method constructs a boundary-expanded base mask, computes a temporal occupancy map from the mask sequence, selects extra temporal-union pixels using radius $r$ and occupancy threshold $\tau$, and runs ProPainter with the refined masks.
        \vspace{0.5\baselineskip}

        \subsection{Experiments}
        Experiments on 100 AI-Hub clips and 30 DAVIS 2017 val clips compare Original, Boundary-only, Temporal union, Ours r10 t0.10, and Ours r10 t0.15. The proposed method improves residue over Boundary-only in most AI-Hub clips and reduces outside changes relative to Temporal union in nearly all AI-Hub clips and all DAVIS clips.
        \vspace{0.5\baselineskip}
        \end{quote}
    }
}
\thispagestyle{empty}
\end{document}



